% Options for packages loaded elsewhere
\PassOptionsToPackage{unicode}{hyperref}
\PassOptionsToPackage{hyphens}{url}
\PassOptionsToPackage{dvipsnames,svgnames,x11names}{xcolor}
%
\documentclass[
  11pt,
]{article}
\usepackage{amsmath,amssymb}
\usepackage{iftex}
\ifPDFTeX
  \usepackage[T1]{fontenc}
  \usepackage[utf8]{inputenc}
  \usepackage{textcomp} % provide euro and other symbols
\else % if luatex or xetex
  \usepackage{unicode-math} % this also loads fontspec
  \defaultfontfeatures{Scale=MatchLowercase}
  \defaultfontfeatures[\rmfamily]{Ligatures=TeX,Scale=1}
\fi
\usepackage{lmodern}
\ifPDFTeX\else
  % xetex/luatex font selection
  \setmainfont[]{Liberation Serif}
  \setsansfont[]{Liberation Sans}
  \setmonofont[]{DejaVu Sans Mono}
\fi
% Use upquote if available, for straight quotes in verbatim environments
\IfFileExists{upquote.sty}{\usepackage{upquote}}{}
\IfFileExists{microtype.sty}{% use microtype if available
  \usepackage[]{microtype}
  \UseMicrotypeSet[protrusion]{basicmath} % disable protrusion for tt fonts
}{}
\makeatletter
\@ifundefined{KOMAClassName}{% if non-KOMA class
  \IfFileExists{parskip.sty}{%
    \usepackage{parskip}
  }{% else
    \setlength{\parindent}{0pt}
    \setlength{\parskip}{6pt plus 2pt minus 1pt}}
}{% if KOMA class
  \KOMAoptions{parskip=half}}
\makeatother
\usepackage{xcolor}
\usepackage[margin=0.8in]{geometry}
\usepackage{longtable,booktabs,array}
\usepackage{calc} % for calculating minipage widths
% Correct order of tables after \paragraph or \subparagraph
\usepackage{etoolbox}
\makeatletter
\patchcmd\longtable{\par}{\if@noskipsec\mbox{}\fi\par}{}{}
\makeatother
% Allow footnotes in longtable head/foot
\IfFileExists{footnotehyper.sty}{\usepackage{footnotehyper}}{\usepackage{footnote}}
\makesavenoteenv{longtable}
\setlength{\emergencystretch}{3em} % prevent overfull lines
\providecommand{\tightlist}{%
  \setlength{\itemsep}{0pt}\setlength{\parskip}{0pt}}
\setcounter{secnumdepth}{-\maxdimen} % remove section numbering
\ifLuaTeX
\usepackage[bidi=basic]{babel}
\else
\usepackage[bidi=default]{babel}
\fi
\babelprovide[main,import]{american}
\ifPDFTeX
\else
\babelfont{rm}[]{Liberation Serif}
\fi
% get rid of language-specific shorthands (see #6817):
\let\LanguageShortHands\languageshorthands
\def\languageshorthands#1{}
\usepackage{adjustbox}
\emergencystretch=3em
\AtBeginEnvironment{longtable}{\small\sloppy}
\usepackage{stmaryrd}
% Printing shims for run-on conditional-bar tokens in the supplied Biology source.
\providecommand{\midL}{\mid L}
\providecommand{\midN}{\mid N}
\sloppy
\ifLuaTeX
  \usepackage{selnolig}  % disable illegal ligatures
\fi
\usepackage{bookmark}
\IfFileExists{xurl.sty}{\usepackage{xurl}}{} % add URL line breaks if available
\urlstyle{same}
\hypersetup{
  pdflang={en-US},
  colorlinks=true,
  linkcolor={blue},
  filecolor={Maroon},
  citecolor={Blue},
  urlcolor={blue},
  pdfcreator={LaTeX via pandoc}}

\author{}
\date{}

\begin{document}

\section{DENYING KNOWLEDGE TO MAKE ROOM FOR
COMMITMENT}\label{denying-knowledge-to-make-room-for-commitment}

\subsection{Truth, Commitment, and Nonmonotonic
Warrant}\label{truth-commitment-and-nonmonotonic-warrant}

\textbf{Gaura Kiśora Dās Rader}\\
\textbf{Heaven≡Earth Foundation}\\
\textbf{Philosophical Companion to the White Paper Series}\\
\textbf{Version 0.10}\\
\textbf{October 2026}

\textbf{Release status.} This manuscript is part of the complete Version
0.10 proposal. Numbered principles identify the framework's current
formulations; their claim status, external validation, and authorization
for consequential application remain separate questions.

\textbf{Companion status.} Active Version 0.10 philosophical companion.
M\&E and \emph{Foundations} control shared definitions, formal scope,
and registered principle wording. This paper controls the extended
philosophical exposition, worked cases, and reproducibility record of
the commitment-logic clarification when consistent with those sources.
This release edition is synchronized to the frozen Version 0.10 corpus
and adds no authority beyond the scope stated here.

\begin{center}\rule{0.5\linewidth}{0.5pt}\end{center}

\section{Abstract}\label{abstract}

Religious and metaphysical discourse has often been forced into a choice
between two unsatisfactory positions. Either its central utterances are
treated as descriptions of hidden states of affairs and judged by
standards they frequently cannot meet, or they are treated as
expressions, practices, or forms of life that no longer answer to
truth---and then appear to answer to nothing at all. The first route
turns faith into speculative science. The second produces the insulation
objection: if commitment is removed from descriptive evaluation without
receiving another standard, nothing internal to the practice can count
against it.

This paper distinguishes a \textbf{description register} from a
\textbf{commitment register} primarily by speech-act function and
standard of assessment, not by positing two exhaustive semantic species.
The same stable content may be asserted, hoped, vowed, trusted, refused,
ritualized, or adopted as a plan. Descriptive uses remain answerable to
truth and evidence. Commitment uses are assessed by warrant, integrity,
protected viability, affected-party standing, consequence, correction,
and repair. The registers are asymmetrically permeable: established
facts constrain commitment, while commitment does not manufacture fact.

For the representable fragment, the paper assumes a Boolean-closed
language whose commitment contents possess stable \textbf{realization
conditions}. It does not derive those conditions and does not offer a
general solution to the Frege--Geach problem. Preferential consequence
operates at a later level: given stable contents, it asks what follows
defeasibly across the best viable realizations of a fixed evaluation
model. Under standard smoothness and definability conditions, the plain
fixed-base preferential consequence relation has the KLM System P
structure. Strong warrant adds feasibility: it coincides with that
relation on feasible premises and is reported separately from
infeasibility. Ranked models support Rational Monotony when relevant
minima are attained.

The resulting practical taxonomy separates \textbf{strong warrant},
\textbf{mere admissibility}, \textbf{exclusion}, \textbf{infeasibility},
and \textbf{underdetermination}. A separate practical-expansion review
asks whether a commitment may responsibly be adopted once facts,
protected floors, affected agents, option-set adequacy, uncertainty,
horizon, reversibility, independent challenge, and revision conditions
are included. The paper's distinctive proposal is therefore not a new
general semantics of commitment. It is the combination of stable
commitment contents with fixed-base nonmonotonic consequence, yielding
distinctions between what every best realization supports and what at
least one supports.

The paper applies this architecture to the insulation objection,
embedded inference, disagreement, transformative commitment, deferred
inquiry, political identity, and a schematic case of two widows who
share the same doctrine but inhabit it differently. It also treats the
register discipline reflexively: adopting it is itself a defeasible
commitment grounded in a broader commitment to truth. Description
remains corrigible by evidence; commitment remains corrigible by fact,
consequence, affected-party testimony, and the trajectories it
generates.

The result is neither fideism nor a general non-cognitivism. It is an
account of how responsible commitment remains possible after counterfeit
knowledge has been denied.

\textbf{Keywords:} commitment; religious language; non-monotonic logic;
preferential semantics; System P; speech acts; Wittgenstein;
Frege--Geach; self-deception; spiritual practice; warrant.

\subsection{Note on Version 0.10}\label{note-on-version-0.10}

This publication preserves the original philosophical project while
incorporating two successive audit cycles. The first separated skeptical
consequence from credulous admissibility and fixed the viability base
within every model to which KLM results are applied. The second
clarified that preferential consequence presupposes rather than supplies
stable content, distinguished descriptive truth-satisfaction from
commitment realization, split infeasibility from underdetermination,
strengthened the governance conditions on model choice, and relocated
the paper's novelty from a general semantics of religious language to a
nonmonotonic consequence architecture over stable commitment contents.

The paper therefore treats the two registers primarily as differences of
principal speech-act function and assessment standard. Formal results
apply only to the Boolean-representable fragment. Ritual, prayer,
poetry, myth, first-person avowal, and other practice-bound forms may be
meaningful while remaining outside this formalization. The complete
correction history and the independent audit disposition are preserved
in the
\href{../supplements/10_Denying_Knowledge_Technical_Supplement_v0.10.md}{technical
supplement}.

This release synchronizes the companion paper with \emph{Coherence
Universalism: Metaphysics and Epistemology} and \emph{Coherence
Universalism: Foundations}. It adds the official CU-COM identifiers and
cross-references to the restored meaning, love, revision, redemption,
and Omega architecture. The formal results, audit disposition, finite
regressions, and philosophical scope established by the audited prior
version are unchanged.

\subsection{Authority and relation to the
series}\label{authority-and-relation-to-the-series}

This paper is the active Version 0.10 philosophical companion for the
extended commitment-logic argument. It owns the detailed philosophical
exposition, worked cases, finite regressions, and the account of how
commitment remains answerable to fact, consequence, affected-party
testimony, correction, and repair.

Its authority is bounded:

\begin{itemize}
\tightlist
\item
  \emph{Coherence Universalism: Metaphysics and Epistemology} controls
  the shared register definitions, fixed-base formal architecture, and
  framework-level scope;
\item
  \emph{Coherence Universalism: Foundations} controls registered
  principle wording, identifiers, statuses, and authority classes;
\item
  the domain papers control their scientific, clinical, social, ethical,
  and engineering claims;
\item
  this companion controls its fuller examples and philosophical
  development only where they remain consistent with those sources; and
\item
  its technical supplement controls the reproducibility record of the
  finite models and counterexamples.
\end{itemize}

When a conflict appears, the controlling source governs until an
explicit, versioned revision is adopted. The companion may expose a need
for upstream revision, but it does not silently create one.

\begin{center}\rule{0.5\linewidth}{0.5pt}\end{center}

\section{Domain Analytic Frame and CAP
Record}\label{domain-analytic-frame-and-cap-record}

\textbf{Target:} Coherence Universalism Version 0.10 release\\
\textbf{Status:} Active companion synchronized through Wave 6; no
independent principle authority; final release freeze pending\\
\textbf{Use rule:} Unresolved fields are left explicit rather than
completed by analogy.

\subsection{Paper and model identity}\label{paper-and-model-identity}

\begin{longtable}[]{@{}
  >{\raggedright\arraybackslash}p{(\columnwidth - 4\tabcolsep) * \real{0.3333}}
  >{\raggedright\arraybackslash}p{(\columnwidth - 4\tabcolsep) * \real{0.3333}}
  >{\raggedright\arraybackslash}p{(\columnwidth - 4\tabcolsep) * \real{0.3333}}@{}}
\toprule\noalign{}
\begin{minipage}[b]{\linewidth}\raggedright
Field
\end{minipage} & \begin{minipage}[b]{\linewidth}\raggedright
Declaration
\end{minipage} & \begin{minipage}[b]{\linewidth}\raggedright
Status
\end{minipage} \\
\midrule\noalign{}
\endhead
\bottomrule\noalign{}
\endlastfoot
Paper & Denying Knowledge to Make Room for Commitment & Specified \\
Domain owner & Philosophical companion / commitment-logic exposition and
worked cases & Active bounded companion scope \\
Primary model or analysis & Fixed-base preferential evaluation model
over worlds or trajectories & Paper-relative \\
Primary claim status & Philosophical argument, definitions, imported KLM
result, formal propositions, practical governance architecture,
interpretations, and commitments & Mixed statuses remain explicit \\
Review gate & Philosophy of religion, logic, semantics and pragmatics,
ethics, psychology, religious studies, and affected-practice review &
Proportionate to claim and consequential reach \\
Methodological authority & M\&E and Foundations control shared
definitions, formal scope, registry wording, and statuses; this
companion controls extended exposition and regressions within that scope
& Bounded companion authority \\
\end{longtable}

\subsection{Shared Analytic Frame}\label{shared-analytic-frame}

\begin{longtable}[]{@{}
  >{\raggedright\arraybackslash}p{(\columnwidth - 4\tabcolsep) * \real{0.3333}}
  >{\raggedright\arraybackslash}p{(\columnwidth - 4\tabcolsep) * \real{0.3333}}
  >{\raggedright\arraybackslash}p{(\columnwidth - 4\tabcolsep) * \real{0.3333}}@{}}
\toprule\noalign{}
\begin{minipage}[b]{\linewidth}\raggedright
\#
\end{minipage} & \begin{minipage}[b]{\linewidth}\raggedright
Slot and domain declaration
\end{minipage} & \begin{minipage}[b]{\linewidth}\raggedright
Status and evidence
\end{minipage} \\
\midrule\noalign{}
\endhead
\bottomrule\noalign{}
\endlastfoot
1 & \textbf{Unit and boundary.} A declared commitment evaluation model
with stable representable contents, descriptive background, protected
constraints, affected-agent scope, horizon, and ordering. & Formal
fragment only \\
2 & \textbf{State or trajectory space.} \(W\) contains worlds, states,
or trajectories; \(V_F\) is the fixed viability base and \(B(F,\Gamma)\)
the minimal best set. & Formally specified \\
3 & \textbf{Environment.} Lived practices, relationships, institutions,
descriptive evidence, and unformalized symbolic or ritual meaning
constrain application from outside the Boolean fragment. & Explicit
boundary \\
4 & \textbf{Constraints and viable region.} \(F,K,A,H\) define the fixed
viability base; facts, floors, affected parties, and horizon are not
generated by KLM structure. & Model-relative \\
5 & \textbf{Actions, options, and transitions.} Premise addition,
commitment adoption, suspension, exclusion, practical expansion,
trajectory choice, and model revision are distinct operations. &
Specified \\
6 & \textbf{Dynamics or update rule.} Preferential filtering is an
inferential operation, not a physical dynamics. Trajectory lifting is
permitted under fixed model parameters. & Formal semantics \\
7 & \textbf{Observation and access.} Evidence, experience, testimony,
records, consequence, and practice determine descriptive facts and
practical review; the formal relation does not observe the world. &
Externally supplied \\
8 & \textbf{Perturbation and intervention class.} New premises within a
fixed model differ from changed facts, floors, affected scope, horizon,
realization conditions, or ordering. & Critical distinction \\
9 & \textbf{Equilibria, basins, or regimes.} Minimal elements under
\(\prec\) are preferential optima, not dynamical equilibria or
psychological attractors. & Formally specified / dynamical status not
applicable \\
10 & \textbf{Local--global status.} Practical expansion examines
affected agents and externalized costs qualitatively; no exact
local--global theorem governs commitments as a whole. & Qualitative \\
11 & \textbf{Scalar status.} The preference relation need not be
generated by one scalar. Strong warrant, admissibility, exclusion,
infeasibility, and underdetermination are distinct verdicts. & No
required scalar \\
12 & \textbf{Statistical-image and allocation status.} The core
preferential model is set- and order-based rather than simplex-valued.
Optional probabilistic extensions require separate CU-M14 licensing. &
No core simplex license \\
13 & \textbf{Translation status.} KLM results reach exact structural
transfer only for the stable Boolean-representable fragment under the
required assumptions; lived application requires further bridges. &
Level 3 formally, lower or higher elsewhere \\
14 & \textbf{Open gaps.} General commitment semantics, practice-bound
meaning, model capture, option adequacy, disagreement across models, and
empirical study of trajectories remain unresolved. & Open \\
\end{longtable}

\subsection{Compact Coherence Analysis Protocol
record}\label{compact-coherence-analysis-protocol-record}

\begin{longtable}[]{@{}
  >{\raggedright\arraybackslash}p{(\columnwidth - 4\tabcolsep) * \real{0.3333}}
  >{\raggedright\arraybackslash}p{(\columnwidth - 4\tabcolsep) * \real{0.3333}}
  >{\raggedright\arraybackslash}p{(\columnwidth - 4\tabcolsep) * \real{0.3333}}@{}}
\toprule\noalign{}
\begin{minipage}[b]{\linewidth}\raggedright
CAP module
\end{minipage} & \begin{minipage}[b]{\linewidth}\raggedright
Declaration
\end{minipage} & \begin{minipage}[b]{\linewidth}\raggedright
Boundary
\end{minipage} \\
\midrule\noalign{}
\endhead
\bottomrule\noalign{}
\endlastfoot
System declaration & Declare \(W,F,K,A,H,V_F,\prec,\Vdash_R\), content
fragment, and option or trajectory space. & Section 3 \\
Observation declaration & State which descriptive facts and lived
consequences are established, disputed, omitted, or institutionally
controlled. & Commitment cannot manufacture fact \\
Evaluation declaration & Distinguish formal verdicts from practical
expansion warrant; declare floors, affected parties, options,
uncertainty, and reversibility. & Consequence is not permission \\
Power and authorship declaration & Name who supplied realization
conditions, viability base, ordering, affected scope, and authority to
apply the verdict. & Captured models remain possible \\
Revision declaration & A changed base, floor, affected scope, horizon,
semantics, or ordering is a new model whose warrant must be
re-established. & Non-inheritance \\
\end{longtable}

\subsection{License summary}\label{license-summary}

\begin{longtable}[]{@{}
  >{\raggedright\arraybackslash}p{(\columnwidth - 4\tabcolsep) * \real{0.3333}}
  >{\raggedright\arraybackslash}p{(\columnwidth - 4\tabcolsep) * \real{0.3333}}
  >{\raggedright\arraybackslash}p{(\columnwidth - 4\tabcolsep) * \real{0.3333}}@{}}
\toprule\noalign{}
\begin{minipage}[b]{\linewidth}\raggedright
License
\end{minipage} & \begin{minipage}[b]{\linewidth}\raggedright
Declaration
\end{minipage} & \begin{minipage}[b]{\linewidth}\raggedright
Current status
\end{minipage} \\
\midrule\noalign{}
\endhead
\bottomrule\noalign{}
\endlastfoot
Scalar license & No scalar required for the preference relation; risk or
decision scalarization must be separately declared. & None by default \\
Local--global license & Qualitative governance review of affected scope
and externalization. & No exact theorem \\
Statistical Image License & Possible only for separately declared
probability-threshold or outcome-distribution extensions. & Optional
extension \\
Allocation License & No allocation structure is required by the
commitment-logic core. & Not applicable \\
Refinement and Coarse-Graining License & No general coarse-graining
theorem is invoked by the fixed-base preferential semantics. & Not
applicable to core \\
Fisher-Rao License & Not licensed for the core logic merely because
probabilities may appear in an extension. & Not applicable to core \\
Translation license & Exact KLM inheritance only under fixed-base,
Boolean, definability, smoothness, ranking, and minimum-attainment
conditions as stated. & Bounded Level 3 \\
\end{longtable}

\subsection{Open gaps and review
disposition}\label{open-gaps-and-review-disposition}

\begin{longtable}[]{@{}
  >{\raggedright\arraybackslash}p{(\columnwidth - 4\tabcolsep) * \real{0.3333}}
  >{\raggedright\arraybackslash}p{(\columnwidth - 4\tabcolsep) * \real{0.3333}}
  >{\raggedright\arraybackslash}p{(\columnwidth - 4\tabcolsep) * \real{0.3333}}@{}}
\toprule\noalign{}
\begin{minipage}[b]{\linewidth}\raggedright
Gap
\end{minipage} & \begin{minipage}[b]{\linewidth}\raggedright
Consequence
\end{minipage} & \begin{minipage}[b]{\linewidth}\raggedright
Required work
\end{minipage} \\
\midrule\noalign{}
\endhead
\bottomrule\noalign{}
\endlastfoot
Stable content & The formalism presupposes rather than explains
realization conditions. & Semantic and practice-sensitive work \\
Model capture & A determinate result can rest on unjust floors or
options. & Independent governance and affected-party standing \\
Trajectory evidence & Illustrative cases do not establish
population-level effects. & Empirical and qualitative research \\
\end{longtable}

\begin{quote}
\textbf{Propagation rule.} This card makes the paper's existing objects
and limits visible. It does not strengthen any empirical, mathematical,
normative, metaphysical, or spiritual claim beyond the wording supported
in the manuscript.
\end{quote}

\begin{center}\rule{0.5\linewidth}{0.5pt}\end{center}

\section{1. Introduction: Two Widows}\label{introduction-two-widows}

Consider a schematic case.

Two women lose a husband within a month of each other. Both belong to
the same congregation, and both hold the same commitments: the dead are
with God, death is not the end, grief is not despair.

The first concludes that to grieve would be to doubt, and does not
grieve. She speaks of her husband's homegoing with composure. Within a
year she cannot say his name aloud. She has stopped visiting the grave,
avoids the couple they used to see on Fridays, and has become brittle
with anyone else in the congregation who is mourning. Asked how she is,
she says she is at peace, and she is not lying.

The second concludes that grief is what love does when it loses its
object, and that hope does not cancel loss. She grieves badly for two
years. She keeps saying his name. At the end of it she is someone who
can sit with other people in their losses, and her faith has not
diminished.

The example stipulates rather than clinically diagnoses the relevant
pattern. Quiet, delayed, private, culturally unfamiliar, or
nonexpressive grief is not by itself evidence of suppression. Behavioral
signs alone never establish the inner organization of a bereaved person.
On the stipulated facts, however, one form of commitment has become
insulated from the loss it claims to interpret, while the other has
remained answerable to it.

The two women share a theological picture. The difference lies in how
that picture is inhabited. Nothing in the case settles whether the dead
are with God. Sincerity does not settle the difference either. What
matters is whether the commitment remains in truthful contact with
experience, agency, relationship, and revision.

This paper asks what standard can locate that difference.

\subsection{1.1 The insulation problem}\label{the-insulation-problem}

A familiar response to evidential criticism of religion holds that
religious utterances do not always function like scientific hypotheses.
They may orient a form of life. They may express trust, dependence,
repentance, gratitude, hope, identity, promise, and practical
allegiance. Demanding the evidential standard appropriate to ordinary
empirical description may therefore misidentify what the utterance is
doing.

This response explains much. It explains why religious language persists
despite evidential pressure that would defeat an empirical hypothesis.
It explains why a creed can remain practically central even when
adherents disagree about literal interpretation. It explains why ritual,
prayer, lament, and blessing cannot be reduced to reports about
invisible objects.

But the response creates a debt. If religious commitment is removed from
descriptive standards and no other standard is supplied, nothing can
count against it. A disciplined commitment and a destructive one become
equally protected. A political or spiritual community can reinterpret
every failure as confirmation. The protection of faith becomes
protection from criticism.

That is the insulation objection.

Its force is not that every religious sentence must be a failed
empirical claim. Its force is that meaningful commitment must remain
vulnerable to corruption and revision even when its principal work is
not descriptive.

\subsection{1.2 The historical
challenge}\label{the-historical-challenge}

The problem inherits the shape of several older debates. Ayer's
verificationist challenge supplied one influential form of the demand
that cognitively meaningful assertion answer to publicly assessable
conditions (Ayer, 1936).

Hume's discussion of miracle testimony asks how extraordinary religious
claims should be weighed against established experience (Hume,
1748/2000, §10). Clifford turns evidential responsibility into a moral
demand: belief formed without adequate evidence can itself be culpable
(Clifford, 1877). Wisdom's invisible-gardener parable, later used in the
Flew--Hare--Mitchell symposium, sharpens the challenge of endlessly
qualified religious assertion (Wisdom, 1945; Flew, 1955; Hare, 1955;
Mitchell, 1955). Geach's challenge to non-cognitivist and ascriptivist
accounts asks how a content can retain its identity when embedded rather
than asserted (Geach, 1960).

One family of replies meets these challenges inside the description
register. Natural theology, Alston's doxastic-practice account, and
Plantinga's proper-function account defend religious belief as truth-apt
and potentially warranted (Alston, 1991; Plantinga, 2000). Nothing in
this paper refutes that project. If a religious proposition purports to
describe reality, it remains answerable to the evidence and arguments
appropriate to that claim.

A second family emphasizes that religious discourse also performs
practical and expressive work. Kant limits theoretical knowledge to make
room for practical faith. Hare's \emph{blik}, Wittgensteinian philosophy
of religion, Phillips's grammar of religious life, and expressivist
approaches treat some utterances as orienting practice rather than
merely picturing hidden fact.

The second family sees something important and leaves something
unresolved. Protecting a practice from a category mistake does not yet
explain what distinguishes a truthful, revisable commitment from an
insane, coercive, self-deceptive, or destructive one.

This paper takes that unresolved problem as its point of departure.

\subsection{1.3 The proposal}\label{the-proposal}

Human discourse has at least two principal functions relevant here.

The \textbf{description register} contains uses whose operative role is
to represent how things are. Such uses are answerable to evidence,
truth, valid inference, calibration, explanatory comparison, and proof
relative to explicit assumptions.

The \textbf{commitment register} contains uses whose operative role is
to orient identity, trust, attention, practice, relationship, or action.
Their central question is not exhausted by whether the embedded content
describes a hidden object. It is whether the commitment can be
responsibly undertaken and inhabited.

The distinction is therefore primarily one of speech-act force,
principal function, and assessment standard. It does not require the
claim that every commitment content belongs to a different semantic
species. ``God is with me'' may function as a literal metaphysical
assertion, a declaration of trust, a ritual avowal, a poetic
condensation of dependence, or a mixed utterance. The proper analysis
decomposes those functions rather than assigning whole traditions to one
category.

The registers are distinct but not sealed:

\begin{quote}
\textbf{Established descriptive facts constrain commitments. Commitments
do not establish descriptive facts.}
\end{quote}

This asymmetry is the first defense against insulation. A commitment
that requires denial of established fact is unwarranted in that respect.
A commitment that helps someone live does not thereby prove its
cosmology.

The second defense is a formal and practical architecture of warrant. It
must distinguish what follows throughout every best realization from
what remains available in only some; what is excluded; what is
infeasible; what the model cannot yet decide; and what may responsibly
be adopted only under wider social and ethical constraints.

\subsection{1.4 The distinctive
contribution}\label{the-distinctive-contribution}

The paper makes six claims.

First, the distinction between description and commitment is best
understood primarily through principal function and assessment standard.
Stable proposition-like content can appear in both registers while being
used and evaluated differently.

Second, commitment-register discourse requires an internal standard.
Sincerity, social acceptance, identity, and psychological usefulness are
insufficient.

Third, the formal apparatus assumes stable realization conditions for a
Boolean-representable fragment. It does not derive a general semantics
for religious, ritual, poetic, or contemplative language, and the KLM
layer is orthogonal to that compositional task.

Fourth, given stable commitment contents, fixed-base preferential
semantics supplies a nonmonotonic consequence relation. It distinguishes
what every best viable realization supports from what at least one
supports.

Fifth, responsible adoption requires more than formal consequence. It
depends on factual compatibility, protected floors, affected-party
standing, option-set integrity, independent challenge, uncertainty,
horizon, reversibility, and revision conditions.

Sixth, adopting the two-register discipline is itself a commitment: a
revisable undertaking to let truth constrain belief and to assess
commitments by what they make of persons, relationships, and
institutions.

The paper's distinctive proposal is therefore narrower than a general
semantics of commitment and stronger for being so:

\begin{quote}
\textbf{Representable commitment contents can retain stable content
while being governed by defeasible rather than classical practical
consequence, yielding distinct verdicts of strong warrant, mere
admissibility, exclusion, infeasibility, and underdetermination.}
\end{quote}

The aim is not to settle every question in philosophy of religion. It is
to show how commitment can remain disciplined after it has ceased to
impersonate knowledge.

\section{2. Content, Force, and Two
Registers}\label{content-force-and-two-registers}

\subsection{2.1 Content is not identical with speech-act
force}\label{content-is-not-identical-with-speech-act-force}

A sentence can contribute stable content while being used in different
ways. ``I will care for my mother'' can be a prediction, a promise, a
vow, a reassurance, or a report of intention. The content relevant to
inference may remain stable while the act performed and the standards
governing it change.

This paper therefore does not define the two registers by a sharp
partition between proposition-like and non-proposition-like meanings. It
distinguishes principal functions:

\begin{longtable}[]{@{}
  >{\raggedright\arraybackslash}p{(\columnwidth - 4\tabcolsep) * \real{0.3333}}
  >{\raggedright\arraybackslash}p{(\columnwidth - 4\tabcolsep) * \real{0.3333}}
  >{\raggedright\arraybackslash}p{(\columnwidth - 4\tabcolsep) * \real{0.3333}}@{}}
\toprule\noalign{}
\begin{minipage}[b]{\linewidth}\raggedright
Principal use
\end{minipage} & \begin{minipage}[b]{\linewidth}\raggedright
Content and force
\end{minipage} & \begin{minipage}[b]{\linewidth}\raggedright
Primary assessment
\end{minipage} \\
\midrule\noalign{}
\endhead
\bottomrule\noalign{}
\endlastfoot
Descriptive assertion & A content is presented as representing how
things are & Evidence, truth, inference \\
Commitment act & A content is adopted, promised, trusted, refused,
consecrated, or used to orient action & Warrant, integrity, trajectory,
agency, correction \\
Symbolic, ritual, poetic, or contemplative performance & Meaning may
depend heavily on practice, embodiment, or first-person use &
Practice-specific interpretation; not automatically covered by the
formal model \\
Semantically empty or manipulatively insulated use & No stable
descriptive, inferential, practical, or corrective role can be
identified & Critique, clarification, or rejection \\
\end{longtable}

The rows are not natural kinds. One utterance may perform several
functions at once.

\subsection{2.2 The description
register}\label{the-description-register}

An utterance functions descriptively when it purports to represent an
independently assessable state of affairs. Examples include:

\begin{itemize}
\tightlist
\item
  a historical claim about an event;
\item
  a causal claim about illness or healing;
\item
  a cosmological claim about the age of the Earth;
\item
  a metaphysical claim used literally;
\item
  a prediction about what a practice will do;
\item
  a psychological claim about grief.
\end{itemize}

Such claims do not leave the description register because they concern
religion, morality, consciousness, or ultimate reality. If they
describe, they answer to descriptive standards.

\subsection{2.3 The commitment register}\label{the-commitment-register}

An utterance functions committally when its principal role is to
organize how an agent or community proceeds. It may vow, bless, mourn,
promise, trust, refuse, consecrate, orient, or receive an event within a
practical horizon.

Commitment-register questions include:

\begin{itemize}
\tightlist
\item
  What does this commitment train the agent to notice or exclude?
\item
  Does it preserve contact with established fact?
\item
  Does it integrate experience or require insulation?
\item
  Does it protect or disable agency?
\item
  Whose costs does it externalize?
\item
  What happens when it encounters grief, conflict, uncertainty, or
  power?
\item
  Can it be revised without destroying the person or community that
  holds it?
\end{itemize}

A commitment act can be unwarranted even when its content is true,
because descriptive truth does not settle whether that content should be
adopted as a vow, plan, identity, or governing orientation.

\subsection{2.4 Mixed discourse}\label{mixed-discourse}

The most important religious utterances are often mixed. Consider:

\begin{quote}
Human beings are expressions of one sacred whole; therefore each person
should be treated with reverence.
\end{quote}

The sentence may contain:

\begin{enumerate}
\def\labelenumi{\arabic{enumi}.}
\tightlist
\item
  a metaphysical assertion or interpretation;
\item
  a commitment to sacred regard;
\item
  a normative conclusion.
\end{enumerate}

The moral attractiveness of the conclusion does not prove the
metaphysical assertion. The commitment need not collapse merely because
the metaphysical picture remains uncertain. The picture can be held as
an interpretation if its status is explicit and it remains compatible
with fact.

Decomposition allows a tradition to preserve meaning without protecting
every inherited factual claim from correction.

\subsection{2.5 How principal function is
attributed}\label{how-principal-function-is-attributed}

Register classification is itself interpretive and can be disputed. No
single grammatical form decides it. Relevant evidence includes:

\begin{itemize}
\tightlist
\item
  what the speaker says they are doing;
\item
  what kinds of correction they accept;
\item
  how the utterance behaves in inference;
\item
  what actions, permissions, or prohibitions it generates;
\item
  how the community treats dissent or failed expectation;
\item
  whether the content is asserted across contexts or used mainly in
  ritual;
\item
  and what institutional consequences depend on taking it literally.
\end{itemize}

Speaker intention matters but is not conclusive. Institutions can use
symbolic language to perform covert descriptive or coercive work.
Conversely, an outside observer can misclassify prayer or lament by
ignoring first-person practice. When the principal function remains
genuinely disputed, the framework should report that uncertainty rather
than force a register assignment.

\subsection{2.6 Asymmetric permeability}\label{asymmetric-permeability}

The relation between registers is one-way with respect to evidential
support.

A commitment is not responsibly available when its operative use
requires systematic denial of established descriptive constraints.

But the converse does not hold. A commitment may be beautiful, healing,
coherent, or morally fruitful without proving its descriptive picture.

Commitments may still influence inquiry by directing attention,
generating hypotheses, sustaining costly investigation, and shaping
which questions are asked. Inquiry is not value-free. Yet motivation and
salience are not evidence, and a commitment's practical warrant does not
establish the truth of the hypotheses it inspires.

This asymmetry blocks two opposed errors:

\begin{itemize}
\tightlist
\item
  \textbf{fideistic insulation:} commitment is treated as beyond factual
  constraint;
\item
  \textbf{reductive elimination:} every form of meaning is forced into
  empirical description.
\end{itemize}

The discipline is to refuse both.

\subsection{2.7 The two-register discipline is itself a
commitment}\label{the-two-register-discipline-is-itself-a-commitment}

The distinction is not a view from nowhere. Its descriptive component is
a philosophical proposal about discourse, inference, and practice.
Choosing to organize inquiry and spiritual life by it is itself a
defeasible commitment.

That commitment is grounded in a broader undertaking:

\begin{quote}
\textbf{Make reality more important than the identity built around one's
interpretation of it.}
\end{quote}

The discipline asks us not to let spiritual desire, fear, belonging, or
moral aspiration determine what is accepted as fact. It also asks us not
to reduce grief, faith, love, purpose, or sacred orientation to failed
science.

Both registers remain revisable. Descriptions change with evidence and
argument. Commitments change when they depend on false premises, require
self-deception, degrade reality contact, foreclose grief or inquiry,
protect domination, or generate trajectories incompatible with truth,
agency, compassion, justice, and repair.

\section{3. The Formal Architecture}\label{the-formal-architecture}

\subsection{3.1 What the formal layer
assumes}\label{what-the-formal-layer-assumes}

The formal layer does not create the meanings on which it operates. It
begins only after a relevant fragment of commitment discourse has been
decomposed into stable, representable contents.

Let:

\begin{itemize}
\tightlist
\item
  \(L_D\) be a descriptive language with ordinary truth-satisfaction
  \(w\vDash_D\delta\);
\item
  \(L_R\) be a declared commitment-representation language;
\item
  \(w\Vdash_R\varphi\) mean that world or trajectory \(w\)
  \textbf{realizes} or \textbf{supports} commitment content \(\varphi\).
\end{itemize}

For the fragment to which the full KLM result is applied, \(L_R\) is
Boolean-closed and the realization extensions satisfy:

\[\adjustbox{max width=\linewidth}{$\displaystyle 
\llbracket\neg\varphi\rrbracket_R
=
W\setminus\llbracket\varphi\rrbracket_R,
$}\]

\[\adjustbox{max width=\linewidth}{$\displaystyle 
\llbracket\varphi\land\psi\rrbracket_R
=
\llbracket\varphi\rrbracket_R
\cap
\llbracket\psi\rrbracket_R,
$}\]

\[\adjustbox{max width=\linewidth}{$\displaystyle 
\llbracket\varphi\lor\psi\rrbracket_R
=
\llbracket\varphi\rrbracket_R
\cup
\llbracket\psi\rrbracket_R.
$}\]

This is an adopted representation of stable commitment content. It is
extension-like and proposition-like enough for embedded inference. A
critic may therefore regard it as relabeling propositions rather than
supplying a new semantics. The paper accepts that limitation. Its
distinctive claim concerns the consequence relation and verdict taxonomy
built over these contents, not the original constitution of their
meaning.

For ritual, prayer, poetry, myth, first-person avowal, and other
contents lacking an adequate Boolean representation, the full System P
claim is not made.

\subsection{3.2 Evaluation models}\label{evaluation-models}

Let:

\[\adjustbox{max width=\linewidth}{$\displaystyle 
M=(W,V_F,\prec,\Vdash_R)
$}\]

be a preferential model for a declared evaluation context.

\begin{itemize}
\tightlist
\item
  \(W\) is a set of worlds, states, or trajectories.
\item
  \(F\) is the established descriptive background.
\item
  \(K\) is the set of protected or hard constraints.
\item
  \(A\) is the affected-agent scope.
\item
  \(H\) is the temporal horizon.
\item
  \(V_F\equiv V_{F,K,A,H}\subseteq W\) is the \textbf{fixed viability
  base} for that model.
\item
  \(\prec\) is a declared preference relation fixed within the model.
\item
  \(\Vdash_R\) supplies the stable realization conditions for the
  representable commitment fragment.
\end{itemize}

For a finite commitment theory \(\Gamma\subseteq L_R\), define:

\[\adjustbox{max width=\linewidth}{$\displaystyle 
\operatorname{Real}_R(\Gamma)
=
\bigcap_{\gamma\in\Gamma}\llbracket\gamma\rrbracket_R,
$}\]

with \(\operatorname{Real}_R(\varnothing)=W\), and:

\[\adjustbox{max width=\linewidth}{$\displaystyle 
W_{F,\Gamma}
=
V_F\cap\operatorname{Real}_R(\Gamma).
$}\]

Premises filter a fixed base by intersection. They do not silently
change protected floors, affected-agent scope, horizon, or preference
order. If any of those changes, the analyst has moved to a different
model. That may be legitimate, but it is model revision rather than
ordinary premise addition.

For \(X\subseteq W\), let:

\[\adjustbox{max width=\linewidth}{$\displaystyle 
\operatorname{Min}_{\prec}(X)
=
\{w\in X:\nexists v\in X\text{ with }v\prec w\}.
$}\]

Write:

\[\adjustbox{max width=\linewidth}{$\displaystyle 
B(F,\Gamma)
=
\operatorname{Min}_{\prec}(W_{F,\Gamma}).
$}\]

The best set may contain one realization, many tied or incomparable
realizations, or---on infinite or non-well-founded structures---none.

The standardized Domain Analytic Frame above is the compact CAP record
for this companion paper. It locates the KLM result at exact structural
transfer only inside the declared formal fragment and prevents that
formal status from silently transferring to model adequacy or practical
permission.

\subsection{3.3 Strong warrant and fixed-base preferential
consequence}\label{strong-warrant-and-fixed-base-preferential-consequence}

Define strong warrant by:

\[\adjustbox{max width=\linewidth}{$\displaystyle 
F,\Gamma\models_C^{\forall}\varphi
\quad\Longleftrightarrow\quad
B(F,\Gamma)\neq\varnothing
\;\land\;
B(F,\Gamma)\subseteq\llbracket\varphi\rrbracket_R.
$}\]

Every best viable realization supports \(\varphi\), and at least one
best realization exists. This is the practical verdict of \textbf{strong
warrant}.

Define the plain fixed-base preferential consequence relation
separately:

\[\adjustbox{max width=\linewidth}{$\displaystyle 
\Gamma\mid\!\sim_C\varphi
\quad\Longleftrightarrow\quad
\operatorname{Min}_{\prec}(V_F\cap\operatorname{Real}_R(\Gamma))
\subseteq
\llbracket\varphi\rrbracket_R.
$}\]

It is not an object-language conditional. It presupposes the realization
semantics of \(L_R\); it does not explain that semantics.

Under the standard definability and smoothness assumptions, this plain
relation satisfies the KLM System P rules (Kraus, Lehmann, \& Magidor,
1990). Ranked extensions support Rational Monotony when the relevant
minima are attained (Lehmann \& Magidor, 1992). Strong warrant coincides
with the plain relation on feasible premises and fails Reflexivity on
infeasible ones, which is why infeasibility is reported as a separate
verdict.

The \textbf{Or} rule is available only because the representable
fragment is Boolean-closed and:

\[\adjustbox{max width=\linewidth}{$\displaystyle 
\llbracket\alpha\lor\beta\rrbracket_R
=
\llbracket\alpha\rrbracket_R
\cup
\llbracket\beta\rrbracket_R.
$}\]

For unformalized sets of practices or commitments without a defined
disjunction, no full System P claim is made.

If the preference is ranked and every relevant nonempty definable set
attains a minimum rank, Rational Monotony is also supported. Total
ordering by itself is not enough when minimum attainment can fail.

These results establish the shape of the inference. They do not supply
the normative or spiritual content of the viability base and ordering.

\subsection{3.4 Admissibility and mere
admissibility}\label{admissibility-and-mere-admissibility}

Define inclusive admissibility:

\[\adjustbox{max width=\linewidth}{$\displaystyle 
F,\Gamma\models_C^{\exists}\varphi
\quad\Longleftrightarrow\quad
B(F,\Gamma)\neq\varnothing
\;\land\;
B(F,\Gamma)\cap\llbracket\varphi\rrbracket_R\neq\varnothing.
$}\]

At least one best viable realization supports \(\varphi\).

A commitment is \textbf{merely admissible} when it is admissible but not
strongly warranted:

\[\adjustbox{max width=\linewidth}{$\displaystyle 
F,\Gamma\models_C^{\exists}\varphi
\quad\land\quad
F,\Gamma\not\models_C^{\forall}\varphi.
$}\]

The distinction matters only where the best set contains more than one
relevant realization. In the three-world finite regression, all 294
merely-admissible verdicts occur in the explicitly isolated
non-singleton-best-set stratum. The frequencies protect the
implementation; they are not evidence that one practical verdict is
common in lived commitment.

Admissibility is weaker than consequence. Two commitments can each be
admissible while their conjunction is not.

Let:

\[\adjustbox{max width=\linewidth}{$\displaystyle 
B(F,\Gamma)=\{w_1,w_2\},
$}\]

with:

\[\adjustbox{max width=\linewidth}{$\displaystyle 
w_1\Vdash_R\varphi\land\neg\psi,
\qquad
w_2\Vdash_R\neg\varphi\land\psi.
$}\]

Then \(\varphi\) and \(\psi\) are each admissible, but
\(\varphi\land\psi\) is not. Admissibility therefore fails \textbf{And}
and is not itself a System P consequence relation.

\subsection{3.5 Exclusion and classical
negation}\label{exclusion-and-classical-negation}

Define:

\[\adjustbox{max width=\linewidth}{$\displaystyle 
F,\Gamma\models_C^{\bot}\varphi
\quad\Longleftrightarrow\quad
B(F,\Gamma)\neq\varnothing
\;\land\;
B(F,\Gamma)\cap\llbracket\varphi\rrbracket_R=\varnothing.
$}\]

No best viable realization supports \(\varphi\). The commitment is
\textbf{excluded} relative to the declared model.

Under the adopted classical complement clause for \(\neg\varphi\),
exclusion is equivalent to strong warrant for the formal negation:

\[\adjustbox{max width=\linewidth}{$\displaystyle 
F,\Gamma\models_C^{\bot}\varphi
\quad\Longleftrightarrow\quad
F,\Gamma\models_C^{\forall}\neg\varphi.
$}\]

This is a simplifying choice with a practical consequence. ``No best
realization supports continuing this practice'' may not be identical in
ordinary deliberation to ``every best realization supports ceasing it.''
The difference can involve non-applicability, abstention, suspension, or
absence of the question. Capturing those distinctions requires a richer
action language or non-classical negation. The present paper does not
undertake that extension.

Exclusion does not mean physical impossibility. It means that
\(\varphi\) is unsupported throughout the best set of the declared
model.

\subsection{3.6 Infeasibility}\label{infeasibility}

The commitment theory is \textbf{infeasible} within the declared model
when:

\[\adjustbox{max width=\linewidth}{$\displaystyle 
W_{F,\Gamma}
=
V_F\cap\operatorname{Real}_R(\Gamma)
=
\varnothing.
$}\]

Nothing survives the conjunction of descriptive background, protected
conditions, and commitment premises. This is different from uncertainty.
It says that the commitment set cannot be realized inside the model as
declared.

\subsection{3.7 Underdetermination}\label{underdetermination}

The case is \textbf{underdetermined} when premise-compatible viable
realizations remain but the evaluation cannot responsibly return one of
the three best-set verdicts. This may occur because:

\begin{itemize}
\tightlist
\item
  a relevant nonempty set has no attained minimum;
\item
  facts remain unsettled;
\item
  affected parties or protected floors are incompletely represented;
\item
  the horizon is unspecified;
\item
  the causal model is inadequate;
\item
  or uncertainty is too great for responsible comparison.
\end{itemize}

Ties do not by themselves create underdetermination. When a nonempty
best set exists, ties are exactly what make mere admissibility
distinguishable from strong warrant and exclusion.

In the finite strict-partial-order regression, every nonempty candidate
set has at least one minimum. The 4,104 empty-best cases are therefore
\textbf{infeasible}, not underdetermined. Underdetermination belongs to
non-attainment and model-adequacy conditions not represented by that
finite suite.

Underdetermination is not permission by default. It is a demand for
inquiry, caution, reversible action, model comparison, or honest
unresolved disagreement.

\subsection{3.8 The inherited maximum-preservation
clause}\label{the-inherited-maximum-preservation-clause}

An earlier formulation asked whether restricting the feasible trajectory
set by adding \(\varphi\) preserved or increased the maximum achievable
score. But the restricted set is a subset of the original set, so its
maximum cannot increase. Equality holds only when at least one old
optimum remains available.

Thus the maximum-preservation condition captures:

\[\adjustbox{max width=\linewidth}{$\displaystyle 
\operatorname{Argmax}(F,\Gamma)
\cap
\llbracket\varphi\rrbracket_R
\neq\varnothing,
$}\]

which is existential support by an old optimum---admissibility---not
skeptical consequence.

The correction does not destroy the project. It reveals that the
original proposal contained two valid questions and treated them as one.

\subsection{3.9 Practical expansion warrant and captured
models}\label{practical-expansion-warrant-and-captured-models}

Formal consequence does not determine whether an agent may responsibly
add a commitment. A commitment can preserve one best trajectory while:

\begin{itemize}
\tightlist
\item
  violating a protected floor;
\item
  imposing concentrated costs on others;
\item
  obtaining local stability through coercion;
\item
  depending on an artificially truncated option set;
\item
  externalizing harm across scales or generations;
\item
  or becoming safe only under an unrealistic model.
\end{itemize}

Let practical expansion warrant be written schematically as:

\[\adjustbox{max width=\linewidth}{$\displaystyle 
\mathcal W_C(\varphi\mid F,\Gamma,K,H,A,M,U).
$}\]

A commitment is responsibly adoptable only if, at minimum:

\begin{enumerate}
\def\labelenumi{\arabic{enumi}.}
\tightlist
\item
  its operative use is compatible with established fact;
\item
  at least one viable continuation remains;
\item
  no inviolable floor is violated;
\item
  affected parties and distributions are represented;
\item
  local integration is not purchased by disabling another agent's
  agency;
\item
  externalized costs across scales and horizons are disclosed;
\item
  the option set has not been controlled so as to predetermine the
  verdict;
\item
  the result is open to sufficiently independent challenge;
\item
  practical correction, appeal, and remedy exist where the commitment
  acquires institutional power;
\item
  the result is sufficiently robust under reasonable model uncertainty;
\item
  the commitment remains revisable;
\item
  where severe irreversible harm is possible, retained option value or a
  necessity-and-proportionality exception is supplied.
\end{enumerate}

Public specification is necessary but insufficient. It makes capture
visible; it does not make the model contestable. The formal apparatus
supplies a verdict structure and no protection against a captured
viability base, ordering, option set, or affected-party model. Those
protections come from the normative governance architecture developed in
\emph{Foundations}, \emph{Social Dynamics}, and \emph{Ethics}.

The verdicts are therefore \textbf{model-relative}. When several people
genuinely share the same model, they can reproduce the same result. The
formalism does not guarantee convergence across rival models, and
transparent unresolved disagreement may be the correct endpoint.

Practical expansion warrant is not a single theorem. It is a governance
architecture for responsible adoption.

\section{4. Addressability and the Boundary of the
Formalization}\label{addressability-and-the-boundary-of-the-formalization}

\subsection{4.1 Meaning exceeds formal
addressability}\label{meaning-exceeds-formal-addressability}

A theory that recognizes commitment-register meaning still needs to
distinguish three questions:

\begin{enumerate}
\def\labelenumi{\arabic{enumi}.}
\tightlist
\item
  Is the utterance meaningful within a practice?
\item
  Is its relevant content stable enough for public criticism and
  embedded inference?
\item
  Is that content representable in the Boolean realization language used
  by the formal model?
\end{enumerate}

The questions do not have the same answer in every case. Prayer, ritual,
poetry, myth, silence, and contemplative avowal may be meaningful
without admitting a complete Boolean decomposition. Such cases are
\textbf{outside the present formalization}, not thereby meaningless.

\subsection{4.2 Addressable commitment
content}\label{addressable-commitment-content}

A commitment is addressable to the degree that its adoption makes a
stable, criticizable difference to what follows. It may:

\begin{itemize}
\tightlist
\item
  require or forbid actions;
\item
  alter which reasons count;
\item
  organize attention;
\item
  change the interpretation of loss or obligation;
\item
  structure relationships;
\item
  create expectations about how the agent should proceed;
\item
  expose the agent to criticism if its own standards are violated;
\item
  and specify conditions of revision or release.
\end{itemize}

For formal treatment, those differences must be representable through
stable realization conditions across the contexts in which the formulas
appear.

\subsection{4.3 A graded criterion}\label{a-graded-criterion}

A commitment \(\varphi\) has commitment-register content to the degree
that adding it to \(\Gamma\) changes:

\begin{itemize}
\tightlist
\item
  the realizable trajectory set;
\item
  the practical priorities or ordering;
\item
  the actions licensed or prohibited;
\item
  the reasons an agent must answer;
\item
  the identity or relationship undertaken;
\item
  or the conditions under which revision is required.
\end{itemize}

It is \textbf{formally addressable} only where enough of that structure
can be represented without falsifying the practice.

This criterion is graded. Some utterances have rich public commitment
content. Others function mainly as ritual, aesthetic form, communal
signal, or first-person orientation. The framework need not force all
meaningful language into one sharply bounded formal class.

\subsection{4.4 Register inheritance}\label{register-inheritance}

A sentence inherits the stricter standard when it performs descriptive
work inside a practical argument.

Suppose a community says:

\begin{quote}
This treatment is divinely required because it always cures the illness.
\end{quote}

The claim about cure is descriptive. Its religious setting does not
exempt it from evidence. If the factual premise fails, the
action-guiding conclusion loses that support even if the community
retains a symbolic or devotional interpretation of the practice.

Register inheritance prevents a mixed utterance from borrowing
insulation from its commitment component while continuing to do
descriptive work.

\subsection{4.5 Not verificationism}\label{not-verificationism}

The account does not say that only empirically testable sentences are
meaningful. It says that different uses answer to different standards,
and that descriptive work cannot be hidden inside a non-descriptive
form.

Mathematics, logic, modality, normativity, ritual, poetry, myth, and
commitment can all possess content. The formal consequence relation
applies only to the addressable fragment for which stable realization
conditions have been responsibly supplied.

\section{5. Frege--Geach and the Boundary of the Formal
Result}\label{fregegeach-and-the-boundary-of-the-formal-result}

\subsection{5.1 The objection}\label{the-objection}

Non-cognitivist and expressivist accounts face a familiar challenge. If
a sentence does not contribute stable content when embedded, how can it
appear in a conditional, negation, or argument without changing meaning?

Consider:

\begin{enumerate}
\def\labelenumi{\arabic{enumi}.}
\tightlist
\item
  If trusting God requires refusing grief, then trusting God can require
  psychological harm.
\item
  Trusting God requires refusing grief.
\item
  Therefore trusting God can require psychological harm.
\end{enumerate}

The embedded and asserted occurrences must contribute sufficiently
stable content for the inference not to equivocate.

\subsection{5.2 What the paper assumes rather than
explains}\label{what-the-paper-assumes-rather-than-explains}

Preferential semantics begins with a language, a realization relation,
and extensions. It presupposes that the formalized contents remain
stable across embedding. KLM consequence does not generate compositional
meaning and cannot, by itself, solve the Frege--Geach problem.

The paper therefore adopts a narrower position.

For the representable fragment, commitment contents are treated as
stable, extension-like, and Boolean-composable. The two registers differ
primarily in speech-act force and assessment standard, not necessarily
in semantic type. A proposition-like content can be asserted as
description, adopted as plan, undertaken as vow, or embedded in
reasoning.

This move avoids claiming that an unstructured attitude becomes
compositional merely because a preferential consequence relation is
placed above it. It also concedes a familiar criticism: assigning
realization extensions can look like relabeling propositions. The paper
does not claim semantic novelty here.

\subsection{5.3 What the KLM layer
contributes}\label{what-the-klm-layer-contributes}

Given stable commitment contents, the formal layer asks a different
question:

\begin{quote}
What follows defeasibly from those contents across the best viable
realizations of a fixed evaluation model?
\end{quote}

The answer distinguishes:

\begin{itemize}
\tightlist
\item
  what every best realization supports;
\item
  what at least one supports;
\item
  what none supports;
\item
  what cannot be realized;
\item
  and what the model cannot yet decide.
\end{itemize}

That is a consequence architecture, not a theory of meaning.

The fixed-base control matters. A countermodel shows that
premise-responsive viability revision can violate Cautious Monotony; the
paired control shows that the same worlds and ordering do not violate it
when premises filter one fixed base. The formal contribution lies in
isolating the source of failure and preserving the distinction between
model revision and inference inside a model.

\subsubsection{5.3.1 Compact countermodel and fixed-base
control}\label{compact-countermodel-and-fixed-base-control}

A three-world construction makes the distinction visible. Let

\[\adjustbox{max width=\linewidth}{$\displaystyle 
W=\{w_1,w_2,w_3\},
$}\]

with \(w_1\vDash b\land g\), \(w_2\vDash\neg b\land g\), and
\(w_3\vDash b\land\neg g\). Order \(w_1\prec w_2\) and \(w_3\prec w_2\),
leaving \(w_1\) and \(w_3\) incomparable. If the viable base is allowed
to change with the premises,

\[\adjustbox{max width=\linewidth}{$\displaystyle 
V_{\varnothing}=\{w_1,w_2\},
\qquad
V_{\{b\}}=\{w_3\},
$}\]

then \(\varnothing\mid\!\sim_C b\) and \(\varnothing\mid\!\sim_C g\),
but \(\{b\}\not\mid\!\sim_C g\). Cautious Monotony fails because adding
\(b\) changes the viable base rather than merely filtering it.

Under the fixed-base control, one declared \(V_F\) is held constant and
premises act only by intersection with
\(\operatorname{Real}_R(\Gamma)\). The same worlds and ordering then
remain within the KLM fixed-base setting. The contrast is not a claim
that model revision is illegitimate; it shows that revision and
inference are different operations and cannot share one theorem without
further proof.

\subsection{5.4 Object-language conditionals and metalinguistic
consequence}\label{object-language-conditionals-and-metalinguistic-consequence}

An object-language conditional belongs to \(L_R\) and receives a
realization extension through the adopted semantic clauses.

The relation:

\[\adjustbox{max width=\linewidth}{$\displaystyle 
\Gamma\mid\!\sim_C\varphi
$}\]

is metalinguistic. It says that \(\varphi\) holds throughout the minimal
\(\Gamma\)-realizing worlds of the fixed viability base. It is not
itself the meaning of ``if,'' ``ought,'' ``promise,'' or any religious
expression.

Conflating the two would make preferential consequence appear to solve a
compositionality problem it presupposes away.

\subsection{5.5 Relation to expressivist and imperative
semantics}\label{relation-to-expressivist-and-imperative-semantics}

Gibbard's planning semantics, Yalcin's expressivist formalism, and work
on imperatives show that practical or directive discourse can possess
sophisticated logical structure without being modeled as ordinary
descriptive assertion (Gibbard, 2003; Yalcin, 2012; Charlow, 2014a,
2014b). They also make the relabeling question unavoidable: if practical
contents are represented by stable sets of worlds or plans, how far do
they differ from propositions?

The present paper takes no general position capable of resolving that
debate. It uses stable realization conditions for a limited practical
purpose and locates its distinctive claim elsewhere: \textbf{classical
or proposition-like content can be governed by nonmonotonic practical
consequence rather than by one undifferentiated classical entailment
relation.}

\subsection{5.6 Scope of the result}\label{scope-of-the-result}

The strongest defensible claim is:

\begin{quote}
\textbf{For commitment contents representable in a declared
Boolean-closed language with stable realization conditions, fixed-base
preferential semantics supplies a defeasible consequence structure that
preserves those contents through embedded inference.}
\end{quote}

This claim concerns consequence, not the original constitution of
meaning.

The result does not cover every use of:

\begin{itemize}
\tightlist
\item
  ritual;
\item
  prayer;
\item
  symbol;
\item
  myth;
\item
  poetry;
\item
  first-person avowal;
\item
  contemplative utterance;
\item
  or embodied religious practice.
\end{itemize}

Those remain outside the formalization unless a responsible
representation is supplied.

The paper therefore does not claim to solve the general Frege--Geach
problem as an expressivist. It assumes stable content for a bounded
fragment and develops a practical consequence architecture over it.

\section{6. Disagreement and Model
Relativity}\label{disagreement-and-model-relativity}

\subsection{6.1 Why disagreement appears
difficult}\label{why-disagreement-appears-difficult}

If commitment uses are not assessed only as ordinary descriptions, what
are people disagreeing about? One answer says that they merely express
different attitudes. But many practical disputes are stronger than
difference. People accuse each other of inconsistency, evasion, cruelty,
corruption, or failure to live by their own stated commitments.

The account distinguishes several cases.

\subsection{6.2 A taxonomy}\label{a-taxonomy}

\subsubsection{Different commitments, correct
verdicts}\label{different-commitments-correct-verdicts}

Two agents may inhabit different commitment sets and face different
trajectory spaces. One commitment may be admissible for one agent and
unavailable to another without either verdict being mistaken.

This is difference, not necessarily disagreement.

\subsubsection{Shared commitments, conflicting
verdicts}\label{shared-commitments-conflicting-verdicts}

If two agents share the same facts, protected floors, affected scope,
horizon, model, and ordering but disagree about whether the formal
condition holds, at least one has misapplied the model.

This is genuine resolvable disagreement.

\subsubsection{Descriptive disagreement}\label{descriptive-disagreement}

Parties may disagree about \(F\): what happened, what a practice causes,
whether testimony is reliable, or whether a risk is real. Ordinary
evidential standards apply.

\subsubsection{Model disagreement}\label{model-disagreement}

Parties may accept the same facts while using different causal,
psychological, or social models. Resolution requires comparative
explanation, prediction, intervention, or acknowledgment of uncertainty.

\subsubsection{Constraint disagreement}\label{constraint-disagreement}

Parties may disagree about which capacities or rights function as
protected floors and which tradeoffs are legitimate. This is a
substantive moral or political disagreement, not a defect that logic
resolves automatically.

\subsubsection{Ordering disagreement}\label{ordering-disagreement}

Parties may rank viable trajectories differently. They must make the
ranking explicit and defend its relation to agency, suffering, truth,
justice, and repair.

\subsubsection{Horizon disagreement}\label{horizon-disagreement}

A commitment can look stabilizing over a week and destructive over five
years. A horizon is not a neutral technical choice. It must be justified
by the phenomenon.

\subsubsection{Affected-agent
disagreement}\label{affected-agent-disagreement}

Parties may omit different people, communities, species, or future
systems. Many apparent disagreements about ``coherence'' are
disagreements about whose trajectory counts.

\subsubsection{Practical conflict}\label{practical-conflict}

Two commitments may be separately warranted yet impossible to realize
together. The result is conflict, negotiation, boundary, or
institutional design---not necessarily an epistemic error.

\subsection{6.3 Determinacy relative to a declared
model}\label{determinacy-relative-to-a-declared-model}

The formal relation is model-relative.

Given the same \(V_F\), \(\prec\), realization semantics, floors,
affected scope, and horizon, different reasoners can reproduce the same
verdict. This is determinate relative to shared public arguments.

It is not a guarantee of intersubjective objectivity across rival
models. Whoever controls the viability base and ordering can control the
verdicts. An institution can omit alternatives, exclude affected people,
or deliberately provide a coarse model that returns no verdict.

Public specification makes those choices visible. It does not make them
contestable.

Contestation requires additional normative architecture:

\begin{itemize}
\tightlist
\item
  affected-party standing;
\item
  option-set integrity;
\item
  independent challenge;
\item
  protected dissent;
\item
  noncompensation for severe concentrated harm;
\item
  practical correctability;
\item
  and authority to revise or reject the model.
\end{itemize}

The formal relation supplies none of these by itself.

\subsection{6.4 Arbitrariness constrained, not disagreement
abolished}\label{arbitrariness-constrained-not-disagreement-abolished}

The account constrains arbitrariness in several ways.

First, descriptive facts remain answerable to evidence regardless of
preference.

Second, stable realization semantics and explicit inference constrain
which conclusions follow inside a model.

Third, protected floors, affected-party scope, and governance
requirements constrain local integration purchased through domination or
exclusion.

Fourth, model revisions must be disclosed rather than smuggled in as
premise addition.

These constraints do not guarantee one final ordering. Cross-model
disagreement may remain transparent and unresolved. Naming that limit is
preferable to pretending the formalism adjudicates values it only
exposes.

\section{7. Asymmetric Permeability}\label{asymmetric-permeability-1}

\subsection{7.1 Description to
commitment}\label{description-to-commitment}

Facts may enter commitment evaluation as constraints.

If a sacred text is shown not to be an accurate geological account, a
commitment that requires treating it as geology loses warrant in that
use. The community may preserve the text as myth, ritual, poetry,
covenant, or theological interpretation, but it may not use the
attractiveness of those functions to re-establish the descriptive claim.

If a healing practice is shown to cause harm, its devotional meaning
cannot exempt it from that evidence.

If a political identity depends on a false historical claim, the
identity does not manufacture the history.

\subsection{7.2 Commitment to
description}\label{commitment-to-description}

The converse evidential path is blocked.

A commitment may motivate inquiry. It may suggest hypotheses, direct
attention, sustain courage, or affect theory choice where evidence
underdetermines alternatives. Testimony and values can reveal neglected
objects of study. None of this makes inquiry value-free.

But practical warrant does not count as evidence that a cosmology or
causal claim is true. Commitments shape inquiry without deciding its
descriptive results.

This distinction is the line the framework commits itself to hold as
disciplinedly as possible.

\subsection{7.3 The answer to the insulation
objection}\label{the-answer-to-the-insulation-objection}

The account does not protect commitment from criticism. It relocates and
pluralizes criticism.

A commitment can fail because it:

\begin{itemize}
\tightlist
\item
  contradicts established fact;
\item
  lacks stable, addressable content for the claim being made;
\item
  violates protected floors;
\item
  requires self-deception;
\item
  disables agency;
\item
  externalizes costs;
\item
  forecloses revision;
\item
  depends on a captured model;
\item
  or produces trajectories rejected throughout the best set of an
  adequately governed evaluation.
\end{itemize}

That is enough to distinguish disciplined from destructive commitment
without pretending to know what cannot presently be known.

\section{8. Transformative Commitment and Model
Revision}\label{transformative-commitment-and-model-revision}

\subsection{8.1 Why instantaneous stability is too
conservative}\label{why-instantaneous-stability-is-too-conservative}

Some warranted transformations begin with disruption.

A person may accept a painful truth that their current identity cannot
yet integrate. A religious practitioner may discover that a foundational
historical claim is false. A patient may receive a diagnosis that
reorganizes their life. A community may admit harm that destabilizes its
public self-understanding.

At the moment of recognition, descriptive adequacy may rise while
integration falls. If warrant requires immediate stability, the
framework will counsel denial precisely when truth demands
reorganization.

\subsection{8.2 Trajectory worlds}\label{trajectory-worlds}

The solution is to evaluate extended trajectories.

A trajectory may contain:

\begin{itemize}
\tightlist
\item
  temporary destabilization;
\item
  grief;
\item
  loss of role;
\item
  a bounded decline in functioning;
\item
  new support;
\item
  reorganization;
\item
  later recovery;
\item
  newly acquired capacities;
\item
  or foreclosed possibilities.
\end{itemize}

The fixed viability predicate evaluates the trajectory as a whole. The
proposition representing adoption selects trajectories in which the
commitment is adopted. It does not redefine the viability base merely by
being added as a premise.

\subsection{8.3 Bounded troughs}\label{bounded-troughs}

A transformative trajectory should be distinguished from romanticized
collapse.

Responsible adoption requires attention to:

\begin{enumerate}
\def\labelenumi{\arabic{enumi}.}
\tightlist
\item
  an inviolable viability floor;
\item
  the depth and duration of the trough;
\item
  credible recovery mechanisms;
\item
  evidence-sensitive expectations;
\item
  available lower-cost alternatives;
\item
  distribution of costs across affected people;
\item
  an expiry condition;
\item
  retained option value where severe harm is possible.
\end{enumerate}

No commitment becomes warranted merely because suffering can be narrated
as growth.

\subsection{8.4 When the model changes}\label{when-the-model-changes}

If an analyst changes protected floors, affected-agent scope, horizon,
causal model, or ordering in response to a commitment, the comparison is
between different models:

\[\adjustbox{max width=\linewidth}{$\displaystyle 
M_0
\longrightarrow
M_1.
$}\]

That can be necessary. Transformative commitments may reveal that the
old model omitted a person, misunderstood a capacity, or used the wrong
horizon.

But the revision must be declared, dated, and justified. A verdict
issued in \(M_0\) does not inherit authority in \(M_1\) without a bridge
argument showing which facts, constraints, and order relations were
preserved. The System P theorem does not transfer automatically to a
non-intersective premise-responsive selection rule.

A commitment must not protect itself by repeatedly redefining harm,
agency, affected parties, or horizon and then presenting each revision
as though it preserved the warrant of the previous model.

\subsection{8.5 Commitment-trajectory
review}\label{commitment-trajectory-review}

A commitment should be reviewed not only by what it says, but by what it
becomes.

Questions include:

\begin{itemize}
\tightlist
\item
  What does it train the person to notice?
\item
  What becomes difficult to acknowledge?
\item
  Which emotions become forbidden?
\item
  Whose testimony becomes unintelligible?
\item
  Which authorities become immune?
\item
  What forms of relationship are produced?
\item
  Can doubt, grief, refusal, and revision still occur?
\item
  What happens when the commitment acquires power?
\end{itemize}

This is the practical discipline of the commitment register.

\section{9. The Two Widows Revisited}\label{the-two-widows-revisited}

\subsection{9.1 Scope and stipulation}\label{scope-and-stipulation}

The example is schematic, conditional, non-clinical, and culturally
non-universal.

Quiet, delayed, private, or nonexpressive grief is not evidence of
pathology or suppression. Behavioral signs alone never establish the
inner organization of bereavement. The case stipulates that avoidance,
brittleness, and loss of relational availability are causally linked to
a rule that grief must be refused as doubt.

The formal question is therefore:

\begin{quote}
Given the stipulated facts and declared model, how are two different
expansions of one theological background evaluated?
\end{quote}

\subsection{9.2 Shared background}\label{shared-background}

Let both women share the commitment set \(\Gamma\):

\begin{itemize}
\tightlist
\item
  the dead are with God;
\item
  death is not the end;
\item
  grief is not despair;
\item
  trust is an appropriate response to what cannot be understood.
\end{itemize}

For the formal comparison, these items enter \(\Gamma\) as commitment
contents: trust, hope, and an orientation toward grief and death. Their
descriptive counterparts are not placed in the established fact set
\(F\); they remain undecided. The paper therefore does not decide
whether the theological descriptions are true. It asks how the same
commitment contents are inhabited, and the comparison is intended to
hold whatever the truth of the undecided descriptions.

Consider two candidate expansions:

\begin{itemize}
\tightlist
\item
  \(\varphi_s\): grief manifests insufficient trust and should be
  suppressed;
\item
  \(\varphi_g\): grief is compatible with hope and may be fully
  undergone.
\end{itemize}

\subsection{9.3 The declared evaluation}\label{the-declared-evaluation}

A responsible model must identify at least:

\begin{itemize}
\tightlist
\item
  the widow herself;
\item
  close relationships and the mourning community;
\item
  affective access and flexibility;
\item
  factual self-contact;
\item
  relational availability;
\item
  ordinary functioning;
\item
  capacity for revision;
\item
  freedom from coercive spiritual bypassing;
\item
  cultural context;
\item
  and a multi-year horizon.
\end{itemize}

These choices are not dictated by logic. They are the substantive human
content of the evaluation.

\subsection{9.4 Suppression}\label{suppression}

Suppression may perform well over a short interval. It can reduce
visible distress, preserve role continuity, and secure communal
approval. It may therefore remain physically possible and locally
stabilizing.

But the stipulated trajectory is different. Grief continues to organize
avoidance outside the avowed self-understanding. The widow cannot say
her husband's name, avoids reminders, and becomes brittle toward others
who mourn.

The case stipulates two insulated organizations:

\begin{itemize}
\tightlist
\item
  an avowed identity in which faith has absorbed the loss;
\item
  an excluded bereavement that continues to shape conduct.
\end{itemize}

Factual self-contact degrades. Relational availability narrows. The
practice becomes less capable of recognizing its own effects.

If every best viable trajectory under the declared model rejects
\(\varphi_s\), then:

\[\adjustbox{max width=\linewidth}{$\displaystyle 
F,\Gamma\models_C^{\bot}\varphi_s.
$}\]

The suppression commitment is excluded without any verdict on the
afterlife.

\subsection{9.5 Grieving with hope}\label{grieving-with-hope}

The second expansion permits the loss to remain loss. Hope does not
cancel grief; it changes the horizon within which grief is borne.

If at least one best viable trajectory supports this form of mourning,
then:

\[\adjustbox{max width=\linewidth}{$\displaystyle 
F,\Gamma\models_C^{\exists}\varphi_g.
$}\]

It is admissible.

If every best viable trajectory supports it relative to the declared
model, then:

\[\adjustbox{max width=\linewidth}{$\displaystyle 
F,\Gamma\models_C^{\forall}\varphi_g.
$}\]

It is strongly warranted.

The story itself does not prove either quantification. It identifies
dimensions a serious model would need to represent.

\subsection{9.6 What the case
establishes}\label{what-the-case-establishes}

The conditional example shows four things.

First, doctrinal identity does not entail practical identity.

Second, sincerity does not guarantee warrant.

Third, a commitment can be criticized without its metaphysical picture
being settled.

Fourth, spiritual hope can remain intact while a harmful interpretation
of that hope is rejected.

Criticism need not choose between accepting an entire doctrinal package
and destroying the practice. Register decomposition permits repair.

\section{10. Deferred Inquiry Under Resource
Constraint}\label{deferred-inquiry-under-resource-constraint}

\subsection{10.1 The difficult case}\label{the-difficult-case}

Suppose an agent encounters credible evidence against a factual claim
central to their religious or ideological community. Their livelihood,
family, status, friendships, and practical identity are embedded in that
community. Immediate acceptance of the evidence may not produce
reorganization. It may produce collapse.

Three responses are available:

\begin{itemize}
\tightlist
\item
  \textbf{rationalize:} absorb the anomaly through an ad hoc
  explanation;
\item
  \textbf{reorganize:} accept the evidence and rebuild the commitment
  structure;
\item
  \textbf{defer:} hold the anomaly open while building the resources
  required for responsible inquiry.
\end{itemize}

\subsection{10.2 Rationalization}\label{rationalization}

Rationalization restores local fit by preventing evidence from altering
protected identity. It changes evidential standards, invents exceptions,
discredits sources selectively, or converts every counterexample into
proof of persecution.

The resulting stability is purchased through degraded adequacy and
foreclosed correction. It is excluded when better viable trajectories
remain available.

\subsection{10.3 Reorganization}\label{reorganization}

Reorganization is often the epistemically direct response. But an agent
may lack the supports required to survive it.

An account that commands immediate detonation of a person's social and
material world may confuse severity with rigor. The fact that a belief
is false does not imply that every route to releasing it is equally
viable.

\subsection{10.4 Deferral}\label{deferral}

Deferral can be warranted when it means:

\begin{itemize}
\tightlist
\item
  not asserting the disputed claim more strongly;
\item
  preserving the anomaly rather than explaining it away;
\item
  seeking independent relationships, advice, and material support;
\item
  building emotional and intellectual capacity;
\item
  preserving a protected route of exit;
\item
  setting an expiry or review condition;
\item
  arranging review outside the community that benefits from continued
  deferral;
\item
  and preparing for later reorganization.
\end{itemize}

Deferral ceases to be warranted when the conditions for inquiry are
present, when capacity-building stops, when the expiry condition is
indefinitely reset without independent review, or when delay continues
to impose serious costs on others.

This is not permission to believe whatever is comforting. The
descriptive claim remains constrained by evidence. The practical
question is how an agent can come into truthful relation with it without
avoidable collapse.

\subsection{10.5 The community's
responsibility}\label{the-communitys-responsibility}

A community that makes exit catastrophic manufactures the conditions
under which deferral becomes rational.

Such a community degrades its own epistemic integrity by:

\begin{itemize}
\tightlist
\item
  raising the cost of registering disconfirmation;
\item
  controlling relationships and resources;
\item
  punishing inquiry;
\item
  isolating members from alternative support;
\item
  converting social dependence into apparent conviction;
\item
  and certifying its own members as perpetually ``not ready'' to
  question.
\end{itemize}

The framework can therefore return paired verdicts:

\begin{itemize}
\tightlist
\item
  deferral may be warranted for the individual;
\item
  the structure making deferral necessary may be unwarranted at the
  collective level.
\end{itemize}

This avoids blaming the agent for responding rationally to constraint
while refusing to launder the institution that created it.

\section{11. Political Identity and the Compassion
Corollary}\label{political-identity-and-the-compassion-corollary}

\subsection{11.1 Why the political case is
harder}\label{why-the-political-case-is-harder}

Religious traditions can sometimes preserve a narrative as myth,
covenant, ritual, or commitment after releasing a factual
interpretation. Political claims often have less room for this move
because they guide collective action through their descriptive content.

Claims about what happened, what a group is like, what a policy does, or
which risks are real are descriptions. Their purpose is to guide action.
Reclassifying them as symbolic meaning may leave them unable to perform
the political work for which they were introduced.

Political identity therefore creates strong pressure toward
rationalization when factual correction threatens community, status,
livelihood, or belonging.

\subsection{11.2 The symmetric
diagnostic}\label{the-symmetric-diagnostic}

The framework does not identify one faction as uniquely vulnerable. The
process question applies symmetrically:

\begin{quote}
Does this community preserve or foreclose the trajectories through which
its members could discover that it is wrong?
\end{quote}

A political community is epistemically degraded when it:

\begin{itemize}
\tightlist
\item
  makes disagreement socially catastrophic;
\item
  supplies instant anomaly-absorbing explanations;
\item
  applies different evidential standards to allies and opponents;
\item
  controls the option set;
\item
  punishes members who update;
\item
  or erases independent records and challenge channels.
\end{itemize}

Symmetry of method does \textbf{not} imply symmetry of evidence. Equal
standards applied to unequal evidence should yield unequal verdicts.
That is the standards working rather than failing.

The truth of a community's present claim does not excuse a structure
incapable of future correction.

\subsection{11.3 Evaluating claims and evaluating
agents}\label{evaluating-claims-and-evaluating-agents}

A false claim and the person holding it require different assessments.

The claim is evaluated against descriptive evidence. The agent is
evaluated relative to their trajectory space, resources, dependencies,
knowledge, control, and available alternatives.

Two agents can hold the same false claim while occupying different
practical positions. One may have a viable route to correction; another
may face catastrophic loss with no support.

This yields the compassion corollary:

\begin{quote}
A person may be responding to a constrained trajectory space rather than
displaying stupidity, dishonesty, or bad character.
\end{quote}

Compassion does not change the truth value of the claim. It changes what
responsible response requires.

The person harmed by a false or oppressive commitment does not owe the
believer education, forgiveness, reconciliation, emotional labor, or
continued relationship. Compassion toward constrained agents must not
become a new burden placed on those who bear the harm.

\subsection{11.4 Changing the trajectory
space}\label{changing-the-trajectory-space}

Where the binding constraint is resource scarcity rather than evidential
access, more argument may address the wrong variable.

Helpful interventions can include:

\begin{itemize}
\tightlist
\item
  relationships outside the community;
\item
  material and legal support;
\item
  protected confidentiality;
\item
  a place to belong after revision;
\item
  slower inquiry;
\item
  trauma-informed support;
\item
  and examples of identity continuity after belief change.
\end{itemize}

\begin{quote}
\textbf{{[}Empirical hypothesis{]}} When identity-relevant revision
carries no viable exit, otherwise sound argument is more likely to
generate resistance, concealment, or defensive entrenchment than when
material, relational, and identity continuity are protected. The claim
requires empirical testing and does not follow from the formal
consequence relation alone.
\end{quote}

\subsection{11.5 Reflexive use}\label{reflexive-use}

This analysis applies reflexively or not at all.

A person who uses the framework only to explain opponents has converted
it into another identity-protective instrument. Its first use should be
examination of the trajectory structures surrounding one's own
commitments.

\section{12. Placement in the Philosophical
Landscape}\label{placement-in-the-philosophical-landscape}

\subsection{12.1 Kant}\label{kant}

The title inherits Kant's declaration that knowledge must be limited to
make room for faith (Kant, 1781/1787/1998, Bxxx). The present proposal
follows the limiting impulse and changes the positive term.

The aim is not to abolish knowledge. It is to deny counterfeit
knowledge: claims of descriptive certainty where the available methods
do not warrant it.

What enters the opened space is not immunity but
commitment---orientation undertaken without pretending that orientation
is evidence.

\subsection{12.2 Wittgenstein}\label{wittgenstein}

The early Wittgenstein supplies a discipline of boundary: where
picturing fails, one should not counterfeit description. The later
Wittgenstein supplies an account of use: meaning is carried by
practices, inferential roles, and forms of life (Wittgenstein,
1921/2001, 1953/2009).

The present account joins these without treating practice as
self-validating. Use helps identify principal function; practical
warrant leaves commitments criticizable.

\subsection{12.3 Truth-apt religious
epistemology}\label{truth-apt-religious-epistemology}

Alston's account of doxastic practices and Plantinga's account of
warranted Christian belief defend the rationality of explicitly
truth-apt religious belief (Alston, 1991; Plantinga, 2000).

The present paper does not displace those projects. It addresses a
different question: how religious and metaphysical commitments remain
criticizable when their principal function in a context is practical,
existential, ritual, or identity-forming rather than the assertion of an
evidentially settled description.

\subsection{12.4 Hare, Phillips, and
Nielsen}\label{hare-phillips-and-nielsen}

Hare's \emph{blik} identifies world-orienting frameworks but lacks a
sufficiently discriminating account of sane and insane bliks (Hare,
1955).

Phillips and Wittgensteinian philosophy of religion preserve the
internal grammar of faith but remain exposed when grammar becomes
insulation (Phillips, 1970).

Nielsen's criticism is therefore decisive: a practice cannot be
protected merely by declaring its standards internal (Nielsen, 1967).

The present account responds by making commitment open to factual
constraint, negative verdict, model comparison, trajectory review, and
external governance.

\subsection{12.5 Carnap and Quine}\label{carnap-and-quine}

Carnap's framework questions illuminate the practical adoption of a
linguistic or conceptual framework (Carnap, 1950). The two-register
architecture is similarly adopted for the work it enables rather than
derived from a neutral language outside all frameworks.

Quine's holism complicates any sharp separation of isolated statements
from wider webs of belief (Quine, 1951). The present account agrees that
commitments and descriptions operate within networks, while retaining a
functional distinction among the kinds of work particular utterances
perform.

\subsection{12.6 Expressivism, plans, imperatives, and
Frege--Geach}\label{expressivism-plans-imperatives-and-fregegeach}

Expressivist and plan-based accounts have developed sophisticated
responses to embedding. Gibbard represents normative judgment through
plans; Yalcin develops Bayesian expressivism; Charlow explores
imperative content and argues that some versions of expressivism can
avoid familiar Frege--Geach problems (Gibbard, 2003; Yalcin, 2012;
Charlow, 2014a, 2014b). Schroeder provides a useful statement of the
historical and contemporary problem (Schroeder, 2008b).

This paper does not claim priority for the idea that non-descriptive or
practical discourse can possess logical structure. Nor does it claim to
provide a new general semantics of that discourse.

Its specific proposal is to assume stable realization conditions for a
bounded, representable commitment fragment and then use fixed-base
nonmonotonic consequence to distinguish:

\begin{itemize}
\tightlist
\item
  support in every best realization;
\item
  support in at least one;
\item
  exclusion from all;
\item
  infeasibility;
\item
  and underdetermination.
\end{itemize}

The novelty, if the proposal is novel, lies in the combination and
practical taxonomy rather than in Boolean semantics or preferential
logic taken separately.

\subsection{12.7 Brandom and
inferentialism}\label{brandom-and-inferentialism}

Inferentialism explains content through the commitments and entitlements
constituted by practices of giving and asking for reasons (Brandom,
1994). The present proposal shares that emphasis and adds
trajectory-level practical evaluation.

The meaning of a commitment includes its inferential role; its
responsible adoption also depends on what kind of life and institution
that role produces.

\subsection{12.8 Price and global
expressivism}\label{price-and-global-expressivism}

Global expressivist approaches emphasize the practical functions
language serves (Price, 2011). The two-register model is less global: it
preserves a robust description register while recognizing commitment
force and other non-descriptive functions.

The asymmetry is deliberate. Commitment does not replace truth; it
begins where truth alone does not settle how to live.

\subsection{12.9 Moral discourse}\label{moral-discourse}

The paper is not a general expressivism about morals.

Moral discourse includes:

\begin{itemize}
\tightlist
\item
  descriptions of harm and capacity;
\item
  normative judgments;
\item
  public commitments;
\item
  institutional authorizations;
\item
  and reasons governing relationship.
\end{itemize}

The Version 0.10 Ethics framework defends substantive principles
concerning standing, agency, welfare, rights, distribution, power,
responsibility, and repair. The two-register architecture can decompose
moral utterances without reducing them all to commitment speech acts.

\section{13. What Commitments This Account
Licenses}\label{what-commitments-this-account-licenses}

\begin{quote}
\textbf{Register note.} ``Licenses'' does not mean ``scientifically
proves'' or ``logically entails.'' The commitments below are
orientations made intelligible by the account. Facts constrain them;
they do not establish facts. They remain answerable to the trajectories
they generate and open to revision.
\end{quote}

\subsection{13.1 CU-COM1 --- Commitment to
Truth}\label{cu-com1-commitment-to-truth}

The first commitment is not to a religious doctrine. It is to truth.

We commit to making reality more important than the identities,
affiliations, hopes, fears, and spiritual interpretations formed around
it.

This means refusing to let what we need to be true decide what we accept
as fact. It also means refusing the opposite error: treating the limits
of knowledge as proof that no responsible faith, metaphysical
interpretation, or sacred orientation is possible.

\subsection{13.2 CU-COM2 --- Discipline of the Two
Registers}\label{cu-com2-discipline-of-the-two-registers}

We commit to holding the distinction between description and commitment
as disciplinedly as possible.

We will not use a commitment's beauty, usefulness, or healing power as
evidence for its descriptive cosmology.

We will not dismiss a commitment merely because it does not function as
an empirical hypothesis.

We will decompose mixed utterances so that each component answers to the
right standard.

\subsection{13.3 CU-COM3 --- Revision in Both
Registers}\label{cu-com3-revision-in-both-registers}

Descriptions remain revisable by evidence and argument.

Commitments remain revisable when they:

\begin{itemize}
\tightlist
\item
  depend on false premises;
\item
  require self-deception;
\item
  degrade reality contact;
\item
  foreclose grief, doubt, or inquiry;
\item
  protect domination;
\item
  erase affected people;
\item
  or generate institutions incapable of correction.
\end{itemize}

No sacred commitment becomes sacred by becoming unrevisable.

\subsection{13.4 The wider commitment
horizon}\label{the-wider-commitment-horizon}

The two-register discipline makes room for substantive commitment; it
does not itself establish which commitments should be adopted. The
Version 0.10 canonical registry now records several commitments that
this paper helps make intellectually possible: \textbf{CU-COM14} on
love, \textbf{CU-COM15} on meaning as discovery and participation,
\textbf{CU-COM16} on freedom of meaning-revision, \textbf{CU-COM17} on
redemptive possibility without justification of evil, \textbf{CU-COM18}
on Heaven on Earth as social vocation, and \textbf{CU-COM19} on enacting
Omega in the present.

Those commitments receive their substantive arguments in the canonical
and domain papers. This companion contributes the discipline by which
they can be held without being disguised as descriptive knowledge: fact
constrains them; their trajectories remain assessable; and none becomes
immune to revision by acquiring sacred language.

\subsection{13.5 Faith after counterfeit
knowledge}\label{faith-after-counterfeit-knowledge}

Faith can be understood as revisable trust: a way of proceeding under
uncertainty without pretending uncertainty has disappeared.

Such faith can orient prayer, love, service, humility, reverence, and
endurance. It can receive a metaphysical picture without converting the
picture into an evidential conclusion.

This is faith with standards rather than faith without risk.

\subsection{13.6 The practice}\label{the-practice}

The two-register architecture is ultimately less a doctrine than a
discipline.

Before asserting, ask: \textbf{Am I describing?} If so, what evidence
would count?

Before committing, ask: \textbf{What does this commitment make possible,
and what does it train me not to see?}

When the two functions are mixed, separate them.

When a commitment generates harm, do not defend it merely because it
once gave meaning.

When evidence destabilizes identity, build the conditions under which
truth can be survived.

When another person cannot yet revise, ask whether the barrier is
argument or trajectory.

And when the framework itself stops serving truth, revise the framework.

\subsection{13.7 A concise vow}\label{a-concise-vow}

\begin{quote}
We commit not to let our commitments decide what we believe as fact.\\
We commit not to use facts as an excuse to abandon meaning, love, or
sacred responsibility.\\
We commit to assess what our commitments make of us and of the
institutions we build.\\
We commit to revise descriptions when evidence requires it and
commitments when truth, consequence, or compassion requires it.\\
We commit to deny counterfeit knowledge in order to make room for
responsible commitment.
\end{quote}

\begin{center}\rule{0.5\linewidth}{0.5pt}\end{center}

\section{14. Conclusion}\label{conclusion}

The philosophy of religion has often treated its central problem as a
contest between evidence and faith. That framing misses an important
possibility.

A stable content can be used descriptively, committally, ritually,
symbolically, or in mixed form. Description remains answerable to truth
and evidence. Commitment remains answerable to fact, practical warrant,
agency, consequence, correction, and repair.

For the representable fragment, the paper assumes stable Boolean
realization conditions. It does not derive them and does not solve the
general Frege--Geach problem. Its formal contribution begins afterward:
fixed-base preferential consequence provides a nonmonotonic structure
for reasoning with those contents.

The resulting architecture distinguishes:

\begin{itemize}
\tightlist
\item
  strong warrant when every best viable realization supports a
  commitment;
\item
  mere admissibility when some but not all do;
\item
  exclusion when none do;
\item
  infeasibility when no premise-compatible viable realization exists;
\item
  underdetermination when the model cannot responsibly decide;
\item
  practical expansion warrant when adoption survives the wider tests of
  fact, floors, affected agents, option sets, independent challenge,
  horizon, uncertainty, reversibility, and revision.
\end{itemize}

The two widows show why those distinctions matter when the relevant
facts are stipulated rather than diagnosed. Identical doctrine can
organize lives differently. One form of hope can become a prohibition on
grief; another can accompany mourning and deepen compassion. The
difference is not captured by truth, evidence, or sincerity alone.

The account does not prove theology. It does not reduce spiritual life
to utility. It does not make coherence a universal scalar, and it does
not ask logic to supply its own values. Public formalization does not
prevent captured models; standing, option-set integrity, correctability,
and noncompensation remain load-bearing normative requirements.

The paper's distinctive proposal is more modest and more defensible than
a general semantics of commitment:

\begin{quote}
\textbf{Stable commitment contents can be governed by defeasible rather
than classical practical consequence, allowing responsible reasoners to
distinguish what every good realization supports from what merely
remains possible in one.}
\end{quote}

Kant denied knowledge to make room for faith. The present proposal
denies counterfeit knowledge to make room for commitment---and then
refuses to let commitment become a refuge from truth.

\section{Appendix A. Formal Results and
Scope}\label{appendix-a.-formal-results-and-scope}

\subsection{A.1 Realization semantics and fixed-base
consequence}\label{a.1-realization-semantics-and-fixed-base-consequence}

Let \(L_R\) be the Boolean-closed representable commitment fragment with
realization relation \(\Vdash_R\). Let:

\[\adjustbox{max width=\linewidth}{$\displaystyle 
M=(W,V_F,\prec,\Vdash_R),
$}\]

with \(V_F\) and \(\prec\) fixed as premises vary. Define:

\[\adjustbox{max width=\linewidth}{$\displaystyle 
\operatorname{Real}_R(\Gamma)
=
\bigcap_{\gamma\in\Gamma}\llbracket\gamma\rrbracket_R,
$}\]

and:

\[\adjustbox{max width=\linewidth}{$\displaystyle 
\Gamma\mid\!\sim_C\varphi
\quad\Longleftrightarrow\quad
\operatorname{Min}_{\prec}
\left(
V_F\cap\operatorname{Real}_R(\Gamma)
\right)
\subseteq
\llbracket\varphi\rrbracket_R.
$}\]

Under standard definability and smoothness conditions, the relation
satisfies System P. The Or rule requires the Boolean clause:

\[\adjustbox{max width=\linewidth}{$\displaystyle 
\llbracket\alpha\lor\beta\rrbracket_R
=
\llbracket\alpha\rrbracket_R
\cup
\llbracket\beta\rrbracket_R.
$}\]

If the preference is ranked and every relevant nonempty definable set
attains a minimum, the relation satisfies Rational Monotony.

The theorem concerns consequence over stable contents. It does not
supply their compositional meaning.

\subsection{A.2 Model-level verdicts}\label{a.2-model-level-verdicts}

Let:

\[\adjustbox{max width=\linewidth}{$\displaystyle 
W_{F,\Gamma}
=
V_F\cap\operatorname{Real}_R(\Gamma),
\qquad
B(F,\Gamma)=\operatorname{Min}_{\prec}(W_{F,\Gamma}).
$}\]

Then:

\begin{itemize}
\tightlist
\item
  \textbf{strong warrant:} \(B\neq\varnothing\) and
  \(B\subseteq\llbracket\varphi\rrbracket_R\);
\item
  \textbf{mere admissibility:} \(B\neq\varnothing\),
  \(B\cap\llbracket\varphi\rrbracket_R\neq\varnothing\), and
  \(B\nsubseteq\llbracket\varphi\rrbracket_R\);
\item
  \textbf{exclusion:} \(B\neq\varnothing\) and
  \(B\cap\llbracket\varphi\rrbracket_R=\varnothing\);
\item
  \textbf{infeasibility:} \(W_{F,\Gamma}=\varnothing\);
\item
  \textbf{underdetermination:} \(W_{F,\Gamma}\neq\varnothing\) but no
  responsible best-set verdict is available because of non-attainment or
  model inadequacy.
\end{itemize}

In finite strict partial orders, every nonempty \(W_{F,\Gamma}\) has a
minimum. The finite regression therefore tests infeasibility but not
non-attainment.

\subsection{A.3 Classical-negation
consequence}\label{a.3-classical-negation-consequence}

With:

\[\adjustbox{max width=\linewidth}{$\displaystyle 
\llbracket\neg\varphi\rrbracket_R
=
W\setminus\llbracket\varphi\rrbracket_R,
$}\]

exclusion of \(\varphi\) is equivalent to strong warrant for
\(\neg\varphi\). This is an adopted simplification. Richer distinctions
among ceasing, suspending, abstaining, and non-applicability require a
richer action language.

\subsection{A.4 Maximum preservation}\label{a.4-maximum-preservation}

Where a real-valued score is used and maxima exist:

\[\adjustbox{max width=\linewidth}{$\displaystyle 
\max_{W_{F,\Gamma\cup\{\varphi\}}}C
\ge
\max_{W_{F,\Gamma}}C
$}\]

can hold only as equality because the first domain is a subset of the
second. Equality is equivalent to support by at least one old maximizer.
It is therefore a test of credulous admissibility, not skeptical
consequence.

\subsection{A.5 Admissibility is not System
P}\label{a.5-admissibility-is-not-system-p}

The two-best-world construction in §3.4 shows that admissibility can
fail conjunction. No System P claim is made for it.

\subsection{A.6 Trajectory lifting}\label{a.6-trajectory-lifting}

Replacing instantaneous worlds with trajectories preserves the abstract
fixed-base preferential construction when:

\begin{itemize}
\tightlist
\item
  the trajectory-level viability base remains fixed;
\item
  the ordering remains fixed;
\item
  and the commitment premise filters trajectories by intersection.
\end{itemize}

Changing protected floors, affected-agent scope, horizon, or ordering is
model revision.

\subsection{A.7 Formal limitation and governance
dependency}\label{a.7-formal-limitation-and-governance-dependency}

The formal results do not determine:

\begin{itemize}
\tightlist
\item
  who counts as affected;
\item
  which floors are inviolable;
\item
  which model is adequate;
\item
  how uncertainty is treated;
\item
  which options should have been considered;
\item
  which ordering is normatively justified;
\item
  or who possesses standing to challenge the model.
\end{itemize}

Public specification gives visibility, not contestability. Responsible
use depends on affected-party standing, option-set integrity,
independent challenge, practical correctability, noncompensation, and
proportional reversibility safeguards.

\section{Appendix B. Revision and Audit
Note}\label{appendix-b.-revision-and-audit-note}

Version 0.10 preserves the original paper's opening case, insulation
objection, register distinction, asymmetric permeability, disagreement
taxonomy, reflexivity, deferred-inquiry analysis, political application,
and compassion corollary.

The first publication pass reconstructed the original formalism by:

\begin{enumerate}
\def\labelenumi{\arabic{enumi}.}
\tightlist
\item
  separating strong warrant, admissibility, exclusion,
  underdetermination, and practical expansion warrant;
\item
  fixing viability and preference within every model claimed to satisfy
  System P or Rational Monotony;
\item
  distinguishing maximum preservation from skeptical consequence;
\item
  representing transformative change inside trajectories;
\item
  retiring a universal scalar coherence score.
\end{enumerate}

The second publication pass incorporated the independent audit by:

\begin{enumerate}
\def\labelenumi{\arabic{enumi}.}
\tightlist
\item
  stating that the KLM layer presupposes stable content and is
  orthogonal to compositional meaning;
\item
  distinguishing descriptive truth-satisfaction \(\vDash_D\) from
  commitment realization \(\Vdash_R\);
\item
  reframing the registers primarily through speech-act function and
  assessment standard;
\item
  narrowing the Frege--Geach claim to stable, Boolean-representable
  commitment contents;
\item
  splitting infeasibility from underdetermination;
\item
  recording the classical-negation duality of exclusion;
\item
  strengthening the dependency on affected-party standing, option-set
  integrity, independent challenge, practical correctability,
  noncompensation, and reversibility;
\item
  making case verdicts visibly conditional on stipulated facts;
\item
  adding a non-singleton-best-set regression stratum;
\item
  normalizing and expanding the bibliography.
\end{enumerate}

The complete proof history, countermodels, regression results,
independent audit, and disposition record appear in the technical
supplement. Wave 6 authority, citation, and cross-reference changes are
recorded in the
\href{../supplements/Framework_Addendum_Companion_Synchronization_v0.10.md}{Companion-Paper
Synchronization Addendum}.

\section{Appendix C. Technical Supplement and
Reproducibility}\label{appendix-c.-technical-supplement-and-reproducibility}

\textbf{Released evidence package.} The current supplement is
\href{../supplements/10_Denying_Knowledge_Technical_Supplement_v0.10.md}{10\_Denying\_Knowledge\_Technical\_Supplement\_v0.10.md}.
Its companion evidence directory is
\href{../evidence/10_Denying_Knowledge/README.md}{10\_Denying\_Knowledge}.
The supplement and evidence directory expose the derivations,
deterministic constructions, scripts, machine-readable results, source
ledger, and audit materials referenced in this manuscript. They do not
elevate a toy model, simulation, or reconstructed check into
target-domain empirical validation.

\section{References}\label{references}

Mitchell, B. (1955). Theology and falsification. In A. Flew \& A.
MacIntyre (Eds.), \emph{New essays in philosophical theology}. SCM
Press.

Alston, W. P. (1991). \emph{Perceiving God: The epistemology of
religious experience}. Cornell University Press.

Ayer, A. J. (1936). \emph{Language, truth and logic}. Gollancz.

Brandom, R. B. (1994). \emph{Making it explicit: Reasoning,
representing, and discursive commitment}. Harvard University Press.

Carnap, R. (1950). Empiricism, semantics, and ontology. \emph{Revue
Internationale de Philosophie, 4}(11), 20--40.

Charlow, N. (2014a). The meaning of imperatives. \emph{Philosophy
Compass, 9}(8), 540--555. \url{https://doi.org/10.1111/phc3.12151}

Charlow, N. (2014b). The problem with the Frege--Geach problem.
\emph{Philosophical Studies, 167}(3), 635--665.
\url{https://doi.org/10.1007/s11098-013-0119-5}

Clifford, W. K. (1877). The ethics of belief. \emph{The Contemporary
Review, 29}, 289--309.

Flew, A. (1955). Theology and falsification. In A. Flew \& A. MacIntyre
(Eds.), \emph{New essays in philosophical theology} (pp.~96--99). SCM
Press.

Geach, P. T. (1960). Ascriptivism. \emph{The Philosophical Review,
69}(2), 221--225. \url{https://doi.org/10.2307/2183506}

Gibbard, A. (2003). \emph{Thinking how to live}. Harvard University
Press.

Hare, R. M. (1955). Theology and falsification. In A. Flew \& A.
MacIntyre (Eds.), \emph{New essays in philosophical theology}
(pp.~99--103). SCM Press.

Hume, D. (2000). \emph{An enquiry concerning human understanding: A
critical edition} (T. L. Beauchamp, Ed.). Clarendon Press. (Original
work published 1748; final authorial edition 1777)

Kant, I. (1998). \emph{Critique of pure reason} (P. Guyer \& A. W. Wood,
Trans.). Cambridge University Press. (Original works published 1781 and
1787)

Kraus, S., Lehmann, D., \& Magidor, M. (1990). Nonmonotonic reasoning,
preferential models and cumulative logics. \emph{Artificial
Intelligence, 44}(1--2), 167--207.
\url{https://doi.org/10.1016/0004-3702(90)90101-5}

Lehmann, D., \& Magidor, M. (1992). What does a conditional knowledge
base entail? \emph{Artificial Intelligence, 55}(1), 1--60.
\url{https://doi.org/10.1016/0004-3702(92)90041-U}

Nielsen, K. (1967). Wittgensteinian fideism. \emph{Philosophy, 42}(161),
191--209. \url{https://doi.org/10.1017/S0031819100001285}

Phillips, D. Z. (1970). \emph{Faith and philosophical enquiry}.
Routledge \& Kegan Paul.

Plantinga, A. (2000). \emph{Warranted Christian belief}. Oxford
University Press.

Price, H. (2011). \emph{Naturalism without mirrors}. Oxford University
Press.

Quine, W. V. O. (1951). Two dogmas of empiricism. \emph{The
Philosophical Review, 60}(1), 20--43.
\url{https://doi.org/10.2307/2181906}

Schroeder, M. (2008a). \emph{Being for: Evaluating the semantic program
of expressivism}. Oxford University Press.

Schroeder, M. (2008b). What is the Frege--Geach problem?
\emph{Philosophy Compass, 3}(4), 703--720.
\url{https://doi.org/10.1111/j.1747-9991.2008.00155.x}

Wisdom, J. (1945). Gods. \emph{Proceedings of the Aristotelian Society,
45}(1), 185--206. \url{https://doi.org/10.1093/aristotelian/45.1.185}

Wittgenstein, L. (2001). \emph{Tractatus logico-philosophicus} (D. F.
Pears \& B. F. McGuinness, Trans.). Routledge. (Original work published
1921)

Wittgenstein, L. (2009). \emph{Philosophical investigations} (P. M. S.
Hacker \& J. Schulte, Eds.; G. E. M. Anscombe, P. M. S. Hacker, \& J.
Schulte, Trans.; 4th ed.). Wiley-Blackwell. (Original work published
1953)

Yalcin, S. (2012). Bayesian expressivism. \emph{Proceedings of the
Aristotelian Society, 112}(2), 123--160.
\url{https://doi.org/10.1111/j.1467-9264.2012.00329.x}

\end{document}
